
\textbf{- Reference:} 
\href{https://www.zama.ai/post/tfhe-deep-dive-part-1}{TFHE Deep Dive: Part I - Ciphertext types}~\cite{tfhe-1}



$ $




\subsection{Setup}

Initially assume $2\le t<q$ and $t\mid q$. The notation $e\xleftarrow{\chi_\sigma}\mathbb Z_q$ means sampling an integer from a specified discrete Gaussian (or a specified rounded Gaussian) and reducing it modulo $q$. Noise bounds concern its integer lift, not the nonnegative residue. A Gaussian has unbounded tails, so correctness is conditional on the bound and parameters must make its failure probability suitably small.
Let $[0, t-1]$ be the plaintext range, and $[0, q-1]$ the ciphertext range, where $t \ll q$ ($t$ is much smaller than $q$). Randomly pick a vector $\vec{s}$ of length $k$ comprising $k$ ternary numbers sampled from $\{-1, 0, 1\}$ as a secret key (denoted as $\vec{s} \xleftarrow{\$} \{-1, 0, 1\}^k$). Let $\Delta = \dfrac{q}{t}$, the scaling factor of plaintext.

\subsection{Encryption}
\label{subsec:lwe-enc}

\begin{enumerate}
\item Suppose we have a plaintext $m \in \mathbb{Z}_t$ to encrypt. \item Randomly pick a vector $\vec{a} \in \mathbb{Z}_q^k$ (of length $k$) as a one-time random public mask (denoted as $\vec{a} \xleftarrow{\$} \mathbb{Z}_q^k$).
\item Randomly pick a small one-time noise $e \in \mathbb{Z}_q$ sampled from the Gaussian distribution $\chi_\sigma$ (denoted as $e \xleftarrow{\chi_\sigma} \mathbb{Z}_q$). 
\item Scale $m$ by $\Delta$, which is to compute $\Delta \cdot m$. This converts $m \in \mathbb{Z}_t$ into $\Delta \cdot m \in \mathbb{Z}_q$.
\item Compute $b = \vec{a} \cdot \vec{s} + \Delta \cdot m + e \in \mathbb{Z}_q$. 
\end{enumerate}

$ $

The LWE encryption formula is summarized as follows:

$ $

\begin{tcolorbox}[title={\textbf{\tboxlabel{\ref*{subsec:lwe-enc}} LWE Encryption}}]
\textbf{\underline{Initial Setup}:} $\Delta = \dfrac{q}{t}$, $\vec{s} \xleftarrow{\$} \{-1, 0, 1\}^k$, where $t$ divides $q$

$ $

$ $

\textbf{\underline{Encryption Input}:} $m \in \mathbb{Z}_t$, $\vec{a} \xleftarrow{\$} \mathbb{Z}_q^k$, $e \xleftarrow{\chi_\sigma} \mathbb{Z}_q$

$ $

%\textbf{Encryption}

\begin{enumerate}
\item Scale up $m \longrightarrow \Delta \cdot m \text{ } \in \mathbb{Z}_q$

\item Compute $b = \vec{a} \cdot \vec{s} + \Delta  m + e \pmod q$
\item $\textsf{LWE}_{\vec{s},\sigma}(\Delta  m + e) = (\vec{a}, b) \text{ } \in \mathbb{Z}_q^{k + 1}$ 
\end{enumerate}

\end{tcolorbox}

\subsection{Decryption}
\label{subsec:lwe-dec}

In the formulas below, the phase is first lifted to a centered representative $[b-\vec a\cdot\vec s]_q$. The notation $\textsf{LWE}^{-1}_{S,\sigma}$ denotes this phase extraction; the subsequent rounding and reduction modulo $t$ complete message decryption. The same convention applies to RLWE and GLWE. Use a fixed nearest-integer rule, such as $\lceil x\rfloor=\lfloor x+1/2\rfloor$.

\begin{enumerate}
\item Given the ciphertext $(\vec{a}, b)$ where $b = \vec{a} \cdot \vec{s} + \Delta \cdot m + e \in \mathbb{Z}_q$, compute $b - \vec{a} \cdot \vec{s}$, which gives the same value as $\Delta \cdot m + e \in \mathbb{Z}_q$. 
\item Round $\Delta \cdot m + e \in \mathbb{Z}_q$ to the nearest multiple of $\Delta$ (i.e., round it as a base $\Delta$ number), which is denoted as $\lceil \Delta \cdot m + e \rfloor_{\Delta}$. This rounding operation successfully eliminates $e$ and gives $\Delta m$, provided $e$ is small enough to be $|e| < \dfrac{\Delta}{2}$. If $|e| \geq \dfrac{\Delta}{2}$, then some of the higher bits of the noise $e$ will overlap with the plaintext $m$, won't be blown away, and will corrupt some lower bits of the plaintext $m$.
\item Compute $\dfrac{\Delta m} {\Delta}$, which is equivalent to right-shifting $\lceil \Delta \cdot m + e \rfloor_{\Delta}$ by $\text{log}_2 \Delta$ bits. (Here we assume $\Delta$ is a power of 2; if $\Delta$ is not a power of 2, scaling up or down $m$ by $\Delta$ is equivalent to multiplying or dividing the value by $\Delta$.) 
\end{enumerate}

$ $

The LWE decryption formula is summarized as follows:

$ $

\begin{tcolorbox}[title={\textbf{\tboxlabel{\ref*{subsec:lwe-dec}} LWE Decryption}}]

\textbf{\underline{Decryption Input}:} $\textsf{ct} = (\vec{a}, b) \text{ } \in \mathbb{Z}_q^{k+1}$

$ $

%\textbf{Decryption}
\begin{enumerate}
\item $\textsf{LWE}^{-1}_{S,\sigma}(\textsf{ct}) = b - \vec{a}\cdot \vec{s} = \Delta  m + e \pmod q$

\item Scale down $\Bigg\lceil\dfrac{ \Delta  m + e } {\Delta}\Bigg\rfloor \bmod t = m \text{ } \in \mathbb{Z}_t$
\end{enumerate}

For correct decryption, the noise $e$ should be $|e| < \dfrac{\Delta}{2}$.

\end{tcolorbox}

During decryption, the secret key owner can subtract $\vec{a} \cdot \vec{s}$ from $b$ because he can directly compute $\vec{a} \cdot \vec{s}$ by using his secret key $\vec{s}$. 

Scaling places neighboring message representatives $\Delta$ apart, so nearest-integer decoding tolerates errors with $|e|<\Delta/2$. If $\Delta$ is a power of two, encryption scaling is a left shift, as in \autoref{fig:scaling}. Decryption must round to the nearest multiple first and only then divide or right-shift. A bare right shift is incorrect for signed errors: $\Delta=8$, $m=1$, and $e=-1$ give $7/8$, whose nearest integer is $1$ but whose truncation is $0$.

\subsubsection{In the Case of $t$ not Dividing $q$}
\label{subsubsec:lwe-noise-bound}

In Summary~\ref*{subsec:lwe-enc} (\autoref{subsec:lwe-enc}), we assumed that $t$ divides $q$. In this case, there is no upper or lower limit on the size of plaintext $m$: its value is allowed to wrap around modulo $t$ indefinitely, yet the decryption works correctly. This is because any wrap-around of $m$
modulo $t$ is carried by the mod-$q$ reduction of the scaled plaintext during decryption.

If $t\nmid q$, set $\Delta=\lfloor q/t\rfloor$ and $r=q-\Delta t$. A plaintext overflow $m=m'+vt$ changes the effective error from $e$ to $e-rv$. Therefore $q\gg t$ alone does not guarantee decryption for arbitrarily many overflows. The derivation below gives an explicit sufficient condition.

For this derivation let $t\ge3$ be odd; primality is unnecessary here, though batching may impose additional conditions (\autoref{subsubsec:poly-vector-transformation-modulus}). Use centered integer representatives: $m'\in[-(t-1)/2,(t-1)/2]$ and ciphertext representatives in $[-q/2,q/2)\cap\mathbb Z$. Write $m=m'+vt$ with integer $v$. Then:

$\left\lfloor\dfrac{q}{t}\right\rfloor\cdot m + e $

$= \left\lfloor\dfrac{q}{t}\right\rfloor\cdot m' + \left\lfloor\dfrac{q}{t}\right\rfloor\cdot vt + e$   \textcolor{red}{ $\rhd$ applying $m = m' + vt$}

$= \left\lfloor\dfrac{q}{t}\right\rfloor\cdot m' + \dfrac{q}{t}\cdot vt - \left(\dfrac{q}{t} - \left\lfloor\dfrac{q}{t}\right\rfloor\right)\cdot vt + e$  

$= \left\lfloor\dfrac{q}{t}\right\rfloor\cdot m' + qv - \left(\dfrac{q}{t} - \left\lfloor\dfrac{q}{t}\right\rfloor\right)\cdot vt + e$

$= \left\lfloor\dfrac{q}{t}\right\rfloor\cdot m' + qv - \epsilon \cdot vt + e$ \textcolor{red}{ $\rhd$ where $\epsilon = \dfrac{q}{t} - \left\lfloor\dfrac{q}{t}\right\rfloor$, a fractional value between $[0, 1)$}

$ $

The actual decoder includes rounding: $\lceil [b-\vec a\cdot\vec s]_q/\Delta\rfloor\bmod t$. Substituting the preceding identity gives:



$\left\lceil\dfrac{1}{\Delta} \cdot \left(b - \vec{a}\cdot\vec{s} \bmod q\right)\right\rfloor \bmod t$

$= \left\lceil\dfrac{1}{\Delta} \cdot \left(\Delta m + e \bmod q\right)\right\rfloor \bmod t$

$= \left\lceil\dfrac{1}{\left\lfloor\frac{q}{t}\right\rfloor} \cdot \left(\left\lfloor\dfrac{q}{t}\right\rfloor \cdot (m' + vt) + e \bmod q\right)\right\rfloor \bmod t$

$= \left\lceil\dfrac{1}{\left\lfloor\frac{q}{t}\right\rfloor} \cdot \left(\left\lfloor\dfrac{q}{t}\right\rfloor \cdot m' + \left\lfloor\dfrac{q}{t}\right\rfloor\cdot vt + e \bmod q\right)\right\rfloor \bmod t$

$= \left\lceil\dfrac{1}{\left\lfloor\frac{q}{t}\right\rfloor} \cdot \left(\left\lfloor\dfrac{q}{t}\right\rfloor \cdot m' + \left(\dfrac{q}{t} - \epsilon \right)\cdot vt + e \bmod q\right)\right\rfloor \bmod t$

$= \left\lceil\dfrac{1}{\left\lfloor\frac{q}{t}\right\rfloor} \cdot \left(\left\lfloor\dfrac{q}{t}\right\rfloor \cdot m' + vq - \epsilon vt + e \bmod q\right)\right\rfloor \bmod t$


$= \left\lceil\dfrac{1}{\left\lfloor\frac{q}{t}\right\rfloor} \cdot \left(\left\lfloor\dfrac{q}{t}\right\rfloor\cdot m'  - \epsilon vt + e \bmod q \right)\right\rfloor \bmod t$ 

$ $

In order for the decryption to work, we ideally want to eliminate the inner modulo $q$ reduction. That is, assuming the centered residue system $\left[-\dfrac{q}{2}, \dfrac{q}{2}\right)\cap\mathbb Z$, we want $-\dfrac{q}{2} \leq \left\lfloor\dfrac{q}{t}\right\rfloor\cdot m'  - \epsilon vt + e < \dfrac{q}{2}$, or more strictly, $\left|\left\lfloor\dfrac{q}{t}\right\rfloor\cdot m'  - \epsilon vt + e\right| < \dfrac{q}{2}$.

$ $

To ensure these conditions hold, we will make a special assumption: $|-\epsilon vt + e| < \dfrac{\Delta}{2}$. Applying this assumption to the expression above, we can derive the following relation:

$\left|\left\lfloor\dfrac{q}{t}\right\rfloor\cdot m'  - \epsilon vt + e\right| $

$ \leq \left|\left\lfloor\dfrac{q}{t}\right\rfloor\cdot m'\right| + \left|- \epsilon vt + e\right|$

$ \leq \left|\left\lfloor\dfrac{q}{t}\right\rfloor\cdot \dfrac{t-1}{2}\right| + \left|- \epsilon vt + e\right|$ \textcolor{red}{ $\rhd$ only odd $t$ is needed for this centered range}

$ \leq \left|\dfrac{q}{t}\cdot \dfrac{t-1}{2}\right| + \left|- \epsilon vt + e\right|$

$ = \dfrac{q}{2} - \dfrac{q}{2t} + \left|- \epsilon vt + e\right|$

$ < \dfrac{q}{2} - \dfrac{q}{2t} + \dfrac{\Delta}{2}$ \textcolor{red}{ $\rhd$ applying our special assumption $|-\epsilon vt + e| < \dfrac{\Delta}{2}$}

$ = \dfrac{q}{2} + \dfrac{1}{2}\cdot \left(\Delta - \dfrac{q}{t}\right)$

$< \dfrac{q}{2}$ \textcolor{red}{ $\rhd$ since $\Delta < \dfrac{q}{t}$}


$ $

Therefore, if we assume $|-\epsilon vt + e| < \dfrac{\Delta}{2}$, then $\left\lfloor\dfrac{q}{t}\right\rfloor\cdot m'  - \epsilon vt + e \bmod q$ can be simplified to $\left\lfloor\dfrac{q}{t}\right\rfloor\cdot m'  - \epsilon vt + e$. We continue with the following derivation: 


$\left\lceil\dfrac{1}{\left\lfloor\frac{q}{t}\right\rfloor} \cdot \left(\left\lfloor\dfrac{q}{t}\right\rfloor\cdot m'  - \epsilon vt + e \bmod q \right)\right\rfloor \bmod t$ 

$= \left\lceil\dfrac{1}{\left\lfloor\frac{q}{t}\right\rfloor} \cdot \left(\left\lfloor\dfrac{q}{t}\right\rfloor\cdot m'  - \epsilon vt + e \right)\right\rfloor \bmod t$ 



$= \left\lceil m' + \dfrac{-\epsilon  vt + e  }{\lfloor\frac{q}{t}\rfloor} \right\rfloor \bmod t$ 

$= m' + \left\lceil \dfrac{-\epsilon vt + e }{\lfloor\frac{q}{t}\rfloor} \right\rfloor \bmod t$ \textcolor{red}{ $\rhd$ since $\lceil m' \rfloor = m'$}

$= m' \bmod t$  \textcolor{red}{ $\rhd$ applying special assumption $|-\epsilon vt + e| < \dfrac{\Delta}{2} = \dfrac{\lfloor\frac{q}{t}\rfloor}{2}$}

Thus $|e-rv|<\Delta/2$, where $r=\epsilon t=q\bmod t$, is a sufficient correctness condition. A simpler conservative bound is $|e|+r|v|<\Delta/2$. It requires a bound on $v$ as well as on the noise. For example, $q=1024$, $t=3$, $\Delta=341$, $m=513=171t$, and $e=0$ give centered phase $-171$ and decoded value $\lceil-171/341\rfloor\bmod3=2$, although $m\bmod3=0$. Increasing $q$ relative only to $t$ and $e$ cannot justify unbounded plaintext overflow.

We can generalize the formula for the plaintext's scaling factor in Summary~\ref*{subsec:lwe-enc} (in \autoref{subsec:lwe-enc}) as $\left\lfloor\dfrac{q}{t}\right\rfloor$, where $t$ is an odd integer.


\begin{tcolorbox}[title={\textbf{\tboxlabel{\ref*{subsubsec:lwe-noise-bound}} Noise Budget for an Odd Plaintext Modulus $\bm{t}$}}]


Given the plaintext's scaling factor $\Delta=\left\lfloor\dfrac{q}{t}\right\rfloor$ and $t$ is an odd integer, the LWE decryption works correctly as long as the error-bounding condition holds:

$|-\epsilon vt + e| < \dfrac{\Delta}{2}$

$ $

, where $\epsilon = \dfrac{q}{t} - \left\lfloor\dfrac{q}{t}\right\rfloor$ is a fractional value between $[0, 1)$, and $v$ accounts for the $t$-overflows of the plaintext $m$.

\end{tcolorbox}

